\section{Conclusion}\label{sec:final-conclusion}

This confirmation report proposes a representation-centric account of
explainability in deep vision models. The central hypothesis is that trained
models develop a transition regime in which generic visual features become
class-relevant spatial geometry. Dataset-wide principal component analysis red--green--blue
(PCA-RGB) visualization provides the initial observation, while train-fitted
bases, a calibration-only selector, trained--random and label-permutation
controls, held-out wrong-prediction analysis, distance-sensitivity tests, and
cross-dataset validation provide the quantitative evidence. The focused
Pre-LN/Post-LN audit shows that the selected transition reorganizes already
decodable class information into a higher-rank, prototype-aligned geometry.
This answers RQ1 within the tested ConvNeXt settings; cross-architecture and
trained-seed replication will determine how broadly the transition location
generalizes.

\methodname{} provides the first methodological use of the observation. Its
distance from adaptive dataset-global robust modes is simple, class-agnostic,
and gradient-free. Completed internal evidence shows competitive object
localization on ResNet-50 and four held-out CNN families after selector
development on ResNet and ConvNeXt. The results also reveal architecture
heterogeneity and a residual advantage for class-conditioned CAMs on some
models. The planned publication study will add the faithfulness, sanity,
semantic, and cross-dataset evidence needed to define the method's explanatory
scope.

\jmethod{} provides a second, distinct test of representation geometry. It
transfers an early relational state, balanced PC1--32 coordinates at selected
transition and late layers, and late affinity flow. At the same time, it
preserves a student-specific bridge and releases spatial constraints near the
end of training.
The audited implementation achieves strong ImageNet-100 results and a
three-seed gain over H9. This is consistent with a benefit from the J1 design
package. However, the contribution of each component remains to be isolated.
J1 is presently a practical peer rather than a demonstrated winner over
ReviewKD.

\textbf{The next stage is to complete and submit the \methodname{} and
\jmethod{} papers.} Transformer work will follow with a broad initial survey;
a focused review and reproducible pilot will then identify the most informative
representation, transition measure, and evaluation target. This staged plan is
feasible because each phase has an explicit evidence gate, publication outcome,
and contingency.
